\chapter{From perimeter tangents to Lipschitz graphs}\label{sec:perimeter-graph-cover}
A tangent measure is a limiting statement at one point. To obtain rectifiability, one must compare nearby points on sets where the relevant scales are uniform. We give that covering argument in codimension one, following the mechanism of \cite{Sim18}, Chapter 3, Theorem 1.6, printed pp.~88--90, as invoked in the proof of Theorem 4.4 on p.~100. The constants below are chosen for transparency, not sharpness. This section proves countable rectifiability; it does not identify the perimeter measure with Hausdorff measure.

\begin{definition}[Mass outside a candidate tangent plane]\label{def:normal-double-cone}
For a unit vector $e\in\mathbb R^n$, $x\in\mathbb R^n$ and $0<\alpha<1$, write
\[
 C_\alpha(x,e)=\{z:|(z-x)\cdot e|\geq\alpha|z-x|\}.
\]
Thus $C_\alpha(x,e)$ is the union of the two closed cones about the normal line. Small mass in $C_\alpha(x,e)\cap B_r(x)$ expresses concentration near the hyperplane $e^\perp$ through $x$. This notation uses the normal component directly; compare the cone notation in \cite{Sim18}, pp.~88--89.
\end{definition}

\begin{lemma}[The reduced boundary and its normal are Borel]\label{lem:reduced-boundary-borel}
For a locally finite-perimeter set $E$ in $U$, the set $\partial^*E\cap U$ is Borel, and $x\mapsto\nu_E(x)$ is Borel on it. In particular it is measurable for every Hausdorff measure. This supplies the measurability needed in the rectifiability statement, starting from \cite{Sim18}, equations (4.2)--(4.3), pp.~96--97.
\end{lemma}
\begin{proof}
Put $\mu=|D\chi_E|$. For each fixed $r>0$, the functions $x\mapsto\mu(B_r(x))$ and the coordinate functions of $x\mapsto D\chi_E(B_r(x))$ are Borel wherever the closed ball lies in $U$. Indeed the kernel $\chi_{\{|y-x|<r\}}$ is Borel in $(x,y)$, and integration against the positive measure, or the positive and negative parts of each derivative component, preserves measurability. Localize to a countable compact exhaustion of $U$ so all relevant integrals are finite. Thus the normalized vector $A_r(x)=D\chi_E(B_r(x))/\mu(B_r(x))$ is Borel on the Borel set where its denominator is positive and finite.

On $U\cap\operatorname{supp}\mu$, every sufficiently small denominator is positive. The existence of a limit of $A_q(x)$ over positive rational $q\downarrow0$ is a countable Cauchy condition; the assertion that the limiting norm is one is also Borel, using $q=1/k$ once those balls lie in $U$. These conditions are equivalent to the full-radius limit in Definition \ref{def:reduced-boundary}. To see the nontrivial implication, for each fixed small $r$ take rational $q_i\uparrow r$. Continuity from below for the finite measure and each signed component gives $A_{q_i}(x)\to A_r(x)$. Every uniform small-radius bound on rational averages therefore passes to the real radius $r$. The Cauchy limit over rational radii is consequently the limit over all radii. Intersecting the stated Borel conditions with the relatively closed support gives the Borel reduced boundary. Its outward normal is the negative pointwise limit of the Borel functions $A_{1/k}$, so is Borel there.
\end{proof}

\begin{lemma}[A pair of points in a normal direction forces cone mass]\label{lem:normal-cone-ball-inclusion}
Let $x\ne y$, put $d=|y-x|$, and let $e,e_x$ be unit vectors with $|e-e_x|\leq1/8$. If
\[
 |(y-x)\cdot e|>\tfrac12d,
\]
then
\[
 B_{d/32}(y)\subset C_{1/4}(x,e_x)\cap B_{2d}(x).
\]
This is the geometric ball inclusion behind the contradiction argument in \cite{Sim18}, pp.~89--90.
\end{lemma}
\begin{proof}
First $|(y-x)\cdot e_x|\geq|(y-x)\cdot e|-|y-x||e-e_x|>3d/8$. For $z\in B_{d/32}(y)$, the triangle inequality gives
\[
 |(z-x)\cdot e_x|>11d/32,\qquad |z-x|<33d/32<2d.
\]
The ratio of the first bound to the second exceeds $1/3$, hence exceeds $1/4$. This proves both inclusions. The absolute values make the estimate valid for either normal direction.
\end{proof}

\begin{proposition}[Uniform density and cone bounds give one Lipschitz graph]\label{prop:uniform-cone-lipschitz-graph}
Let $m=n-1\geq1$, let $\mu$ be a positive Radon measure on an open set $U\subset\mathbb R^n$, and let $S\subset U$. Suppose there are $c>0$, $\rho>0$, a fixed unit vector $e$, and unit vectors $e_x$ for $x\in S$ such that $|e_x-e|\leq1/8$, $B_\rho(x)\subset U$, and for every $x\in S$ and $0<r<\rho$,
\[
 \mu(B_r(x))\geq c r^m,
 \qquad \mu(C_{1/4}(x,e_x)\cap B_r(x))\leq\delta r^m,
 \qquad \delta=\frac{c}{2\cdot64^m}.
\]
If $\operatorname{diam}S<\rho/2$, then $S$ is contained in the graph of a real function on $e^\perp$ with Lipschitz constant at most $1/\sqrt3$. No measurability assumption on $S$ is required for this geometric containment.
\end{proposition}
\begin{proof}
Take distinct $x,y\in S$ and let $d=|y-x|<\rho/2$. If their difference satisfied the hypothesis of Lemma \ref{lem:normal-cone-ball-inclusion}, the mass of the ball from that lemma would be bounded below and above by
\[
 c(d/32)^m
 \leq\mu(B_{d/32}(y))
 \leq\mu(C_{1/4}(x,e_x)\cap B_{2d}(x))
 \leq\delta(2d)^m=\tfrac12c(d/32)^m,
\]
which is impossible. Hence $|(y-x)\cdot e|\leq|y-x|/2$ for all pairs. With $p$ the orthogonal projection onto $e^\perp$, the Pythagorean identity gives
\[
 |p(y)-p(x)|\geq\tfrac{\sqrt3}{2}|y-x|,
 \qquad |(y-x)\cdot e|\leq\tfrac1{\sqrt3}|p(y)-p(x)|.
\]
In particular $p$ is injective on $S$, so the height function $h(p(x))=x\cdot e$ is well defined and $1/\sqrt3$-Lipschitz on $p(S)$. If $S$ is empty there is nothing to prove. Otherwise Proposition \ref{prop:mcshane-extension}, applied in the Euclidean hyperplane $e^\perp$, extends $h$ to the whole hyperplane with the same Lipschitz bound. Its graph contains $S$.
\end{proof}

\begin{theorem}[A pointwise density-and-cone criterion for a countable graph cover]\label{thm:pointwise-cone-graph-cover}
Let $\mu$ be a positive Radon measure on $U\subset\mathbb R^n$, $m=n-1\geq1$, and $S\subset U$. Suppose that for every $x\in S$ there is a unit vector $e_x$ such that
\[
 \liminf_{r\downarrow0}\frac{\mu(B_r(x))}{r^m}>0,
 \qquad
 \lim_{r\downarrow0}\frac{\mu(C_{1/4}(x,e_x)\cap B_r(x))}{r^m}=0.
\]
Then $S$ is contained in a countable union of Lipschitz graphs over $(n-1)$-planes. This is the hypersurface version of the graph-cover mechanism in \cite{Sim18}, pp.~88--90.
\end{theorem}
\begin{proof}
For each positive integer $k$, set $c_k=1/k$ and $\delta_k=c_k/(2\cdot64^m)$. For positive integers $k,j$, let $S_{k,j}$ contain those $x\in S$ for which $B_{1/j}(x)\subset U$ and both
\[
 \mu(B_r(x))\geq c_kr^m,\qquad
 \mu(C_{1/4}(x,e_x)\cap B_r(x))\leq\delta_kr^m
\]
hold for every $0<r<1/j$. The two pointwise limits ensure that every $x\in S$ belongs to at least one such set: first choose $k$ with $1/k$ strictly below the positive lower limit, then $j$ sufficiently large for both bounds and the domain inclusion. Thus $S=\bigcup_{k,j}S_{k,j}$.

Choose finitely many unit vectors whose open $1/8$-balls cover the unit sphere. Split each $S_{k,j}$ according to which of these vectors is within $1/8$ of $e_x$, allowing overlap. Intersect further with a countable cover of Euclidean space by balls of diameter less than $1/(2j)$. Each resulting piece satisfies Proposition \ref{prop:uniform-cone-lipschitz-graph} with $\rho=1/j$ and $c=c_k$, so lies in one whole Lipschitz graph. There are countably many pieces. Their covering property gives the conclusion. This direct partition uses uniform bounds on each piece and does not need to claim the convergence is uniform on all of $S$.
\end{proof}

\begin{theorem}[Countable rectifiability of the reduced boundary]\label{thm:reduced-boundary-graph-cover}
If $E$ has locally finite perimeter in $U\subset\mathbb R^n$, $n\geq2$, then $\partial^*E\cap U$ is contained in a countable union of Lipschitz graphs over $(n-1)$-planes. In particular it is countably $(n-1)$-rectifiable. This proves the graph-cover conclusion in \cite{Sim18}, Chapter 3, Theorem 4.4, p.~100. The equality $|D\chi_E|=\mathcal H^{n-1}\llcorner\partial^*E$ is not used in this proof.
\end{theorem}
\begin{proof}
Take $\mu=|D\chi_E|$, $S=\partial^*E\cap U$, and $e_x=\nu_E(x)$. Corollary \ref{cor:reduced-boundary-exact-densities} gives $\mu(B_r(x))/r^{n-1}\to\omega_{n-1}>0$ at every point of $S$. Corollary \ref{cor:reduced-boundary-cone-concentration} gives the required cone limit for $\alpha=1/4$. Apply Theorem \ref{thm:pointwise-cone-graph-cover}. Lemma \ref{lem:reduced-boundary-borel} supplies Hausdorff measurability. The graph map $z\mapsto z+h(z)e$ is Lipschitz, so these are also countably many Lipschitz images in the sense of Definition \ref{def:rectifiable-set}. The proof covers every reduced-boundary point, rather than introducing an unexamined exceptional subset. Identification of the boundary measure remains a separate step in the full structure theorem.
\end{proof}
